\section{Nuevas relaciones de producción: capacidad del lado de la oferta, redes de conexión y reescritura de la organización}

La sección anterior trató sobre modelos y productos; esta trata sobre las relaciones de producción. Wu Minghui (吴明辉) sostiene que la IA reestructurará la industria del software, los servicios empresariales y el internet industrial, y que generará nuevas relaciones de conexión y de producción. Las más expuestas son las empresas que solo ofrecen una red de conexión, no tienen capacidad del lado de la oferta y cuya red, además, no es un nodo raíz. La razón es que la IA reducirá el costo de conectar y de procesar información, lo que debilitará el valor de la capa intermedia.

\begin{figure}[H]
\centering
\includegraphics[width=0.94\textwidth]{supply-demand-relation.png}
\caption{Capacidad del lado de la oferta frente a red de conexión: una red que solo conecta, sin capacidad del lado de la oferta, queda en una posición de riesgo. Diagrama conceptual propio, elaborado a partir de la conversación de 02:31:36--02:34:28.}
\end{figure}

\begin{knowledgebox}{Lectura de la figura: la IA le pondrá un nuevo precio a la capa intermedia}
Si una empresa solo se dedica a la intermediación, a la distribución o a mantener cadenas de relaciones, pero no tiene capacidad de producto, datos, herramientas, modelos ni entrega, su valor se reducirá en cuanto un Agent pueda conectar directamente la demanda con la oferta. La capacidad del lado de la oferta es la capacidad de producir resultados reales.
\end{knowledgebox}

\subsection{Las nuevas relaciones de producción: personas, modelos, herramientas y datos se reparten el trabajo de otra forma}

La sección anterior explicó cómo se le pone un nuevo precio a la capa intermedia; esta examina cómo se descompone el trabajo mismo. Las nuevas relaciones de producción que trae la IA no consisten en que «la IA sustituya a todos», sino en una nueva división del trabajo entre personas, modelos, herramientas, datos y organizaciones. El modelo se encarga de la cognición, las herramientas aportan las interfaces de ejecución, los datos proporcionan el contexto, el Agent se ocupa de la colaboración automática y las personas asumen los objetivos, el juicio y la responsabilidad. El futuro de los servicios empresariales no se limita a las API: consiste en reorganizar la producción en torno al resultado de cada tarea.

\begin{figure}[H]
\centering
\includegraphics[width=0.92\textwidth]{new-production-relation.png}
\caption{Las nuevas relaciones de producción que trae la IA: modelos, herramientas, Agent, datos privados y socios humanos se reparten el trabajo de otra forma. Diagrama conceptual propio, elaborado a partir de la conversación de 00:00:42--00:00:55 y 02:31:36--02:34:28.}
\end{figure}

\begin{importantbox}{Criterio para reconocer un cambio en las relaciones de producción}
Si la IA solo mejora la eficiencia de un punto aislado, se trata de una mejora de herramientas; si la IA cambia quién posee los datos, quién ejecuta las tareas, quién asume la responsabilidad y quién obtiene los beneficios, se trata de un cambio en las relaciones de producción.
\end{importantbox}

\begin{knowledgebox}{Términos clave: conexión, lado de la oferta y relaciones de producción}
\begin{tabularx}{\textwidth}{p{0.22\textwidth}p{0.32\textwidth}X}
\toprule
Término & Significado en la entrevista & Cambio en la era de la IA \\
\midrule
Red de conexión & Red que conecta clientes, comercios, información o servicios & El Agent reduce el costo de conectar y el valor de la capa intermedia se reduce. \\
Capacidad del lado de la oferta & Capacidad de producir efectivamente bienes, servicios, modelos o resultados entregables & Cuanto más cerca del resultado, más difícil es prescindir de ella. \\
Relaciones de producción & Quién organiza la producción, quién la ejecuta y quién reparte los beneficios y la responsabilidad & Con la incorporación de modelos y Agent, los roles se redistribuyen. \\
Socio & Persona capaz de responder por los resultados, definir problemas y coordinarse con Agent & Disminuyen los puestos de ejecución repetitiva y aumentan los roles de alta responsabilidad. \\
\bottomrule
\end{tabularx}
\end{knowledgebox}

\subsection{Dueños y socios: una visión de la forma organizacional}

En la entrevista aparece también un juicio extremo: en el futuro, las empresas podrían tener solo dos tipos de personas, dueños y socios. Aquí «socio» no designa a un empleado tradicional, sino a alguien que responde por los resultados junto con la IA, los Agent y las herramientas. Puede que este juicio no se cumpla al pie de la letra, pero expresa una tendencia: cuando la ejecución quede automatizada por la IA, la organización dará más peso a la definición de objetivos, la integración de recursos, la asunción de responsabilidades y la colaboración de alta densidad.

\begin{figure}[H]
\centering
\includegraphics[width=0.94\textwidth]{boss-partner-org.png}
\caption{Organización de dueños y socios: la empresa podría pasar de una organización de empleados a una de pocos dueños y muchos socios de IA. Diagrama conceptual propio, elaborado a partir de la conversación de 02:53:20--03:16:40.}
\end{figure}

\begin{warningbox}{«Solo dueños y socios» no significa que la organización desaparezca de inmediato}
Las empresas reales siguen necesitando puestos, procesos, cumplimiento normativo y gestión. Sin embargo, a largo plazo, los puestos de ejecución repetitiva se reducirán por efecto de la IA, y las personas capaces de definir problemas, integrar recursos y responder por los resultados se parecerán cada vez más a socios.
\end{warningbox}

\subsection{Resumen de la sección}

El impacto de la IA en los servicios empresariales no se limita a una mejora de eficiencia. Reescribirá las redes de conexión, la capacidad del lado de la oferta, la estructura organizacional y el reparto de responsabilidades. Las empresas realmente expuestas son las que solo aportan valor como capa intermedia y carecen de capacidad productiva.
